\documentclass[handout]{beamer}
\usetheme{Madrid}

\usepackage[spanish]{babel}
\usepackage{amsthm}

\usepackage{tikz}
\usetikzlibrary{arrows.meta, positioning, shadows.blur, shapes.misc, calc, shapes.geometric}

\newtheorem{teorema}{Teorema}[section]
\newtheorem{definicion}{Definici\'on}[section]
%\useoutertheme{miniframes}
%\useoutertheme{infolines}
%\useoutertheme{smoothbars}
%\useoutertheme{sidebar}
\useoutertheme{split}
%\useoutertheme{shadow}
%\useoutertheme{tree}
%\useoutertheme{smoothtree}
\usepackage[utf8]{inputenc}
\usepackage{cancel}
\usepackage{listings}
\usepackage{amsthm}
\usepackage{xcolor}
\usepackage{wasysym}
\newcommand{\St}{\mathcal{S}}
\newcommand{\Act}{\mathcal{A}}
\newcommand{\diam}{\diamondsuit}
\newcommand{\cuad}{\square}
\newcommand{\G}{\mathcal{G}}
\newcommand{\K}{\mathcal{K}}
\newcommand{\Gk}{\mathcal{G}_\mathcal{K}}
\newcommand{\dist}{\textsf{Dist}}
\newcommand{\M}{\mathcal{M}}
\newcommand{\Pm}{\mathbf{T}}
\newcommand{\alc}{\diamondsuit}
\newcommand{\until}{\mathcal{ U }}
\newcommand{\Vs}{\mathcal{V}_s}
\newcommand{\supp}{\textsf{supp}}
\newcommand{\Outc}{\mathrm{Outc}}
\newcommand{\Inf}{\textsf{Inf}}

\newcounter{saveenumi}
\newcommand{\seti}{\setcounter{saveenumi}{\value{enumi}}}
\newcommand{\conti}{\setcounter{enumi}{\value{saveenumi}}}

\usepackage{pifont}


\newcommand{\siempevent}{\cuad \diam}
\newcommand{\eventsiemp}{\diam \cuad}
\newcommand{\nextltl}{\bigcirc}
\newcommand{\Poly}{\textsf{Poly}}
\newcommand{\DPoly}{\textsf{DPoly}}
\newcommand{\Dist}{\textsf{Dist}}
\newcommand{\Paths}{\textsf{Paths}} %% ver que hay paths en minuscula
\newcommand{\Pathsfin}{\textsf{Paths}_{\textsf{fin}}}
\newcommand{\Salg}{\bm{\Sigma}}
\newcommand{\PMeas}{\mathsf{PMeas}}
\newcommand{\diff}{\mathbf{d}}
\newcommand{\Hk}{\mathcal{H}_\K}
\newcommand{\Prob}{\mathbb{P}}
\newcommand{\W}{\mathcal{W}}
\newcommand{\DGk}{Derand(\Gk)}
\newcommand{\picuad}{\pi_{\square}}
\newcommand{\pidiam}{\pi_{\diamondsuit}}
\newcommand{\ProbG}{\Prob_{\Gk, s}^{\picuad, \pidiam}}

\usepackage{listings}
\lstset{
  language=Haskell,
  basicstyle=\ttfamily,
  keywordstyle=\color{blue},
  commentstyle=\color{green!40!black},
  stringstyle=\color{orange},
  showstringspaces=false,
  breaklines=true,
  frame=single,
  numbers=left,
  numberstyle=\tiny,
  numbersep=5pt,
  xleftmargin=10pt,
}


\makeatletter
\newenvironment{noheadline}{
    \setbeamertemplate{headline}{}
    \addtobeamertemplate{frametitle}{\vspace*{-0.9\baselineskip}}{}
}{}
\makeatother


\author[author1]{author1\inst{1}\\[1ex]  {\small supervisor1\inst{1,2}}}

\title{Juegos estoc\'asticos polit\'opicos\\ con objetivos de Rabin}
\subtitle{}
\author{Katherine Sullivan}
\institute{Facultad de Ciencias Exactas, Ingenier\'ia y Argimensura \\Universidad Nacional de Rosario}
\date{\today}

\begin{document}

\begin{frame}
	\titlepage
\end{frame}

\AtBeginSection[]
{
	\begin{frame}<beamer>
		\frametitle{\'Indice}
		\tableofcontents[currentsection]
	\end{frame}
}

\AtBeginSubsection[]{
	\begin{frame}<beamer>
		\frametitle{\'Indice}
		\tableofcontents[currentsection, currentsubsection]
	\end{frame}
}

\begin{frame}
	\frametitle{\'Indice}
	\tableofcontents
\end{frame}

\section{Introducci\'on}

\subsection{¿Qu\'e son los juegos estoc\'asticos polit\'opicos?}

\begin{frame}{Una primera respuesta vaga}
	\begin{block}{Juego estoc\'astico polit\'opico I}
		Un juego estoc\'asticos polit\'opicos es un \textbf{modelo matem\'atico} que utilizaremos para \textbf{representar un sistema} en el contexto de \textbf{verificaci\'on}.
	\end{block}

	La verificaci\'on es el proceso de demostrar que un \textbf{sistema cumple con
		una especificaci\'on} formalmente definida.

	% Como decimos es un modelo matematico, esto va a requerir entonces dar una def formal de juego estocastico politopico, para eso vamos a dividir esa explicacion en dos partes: primero la de juego estocastico y despues la de politopico

\end{frame}

% \begin{frame}{Caracterizando algunos modelos}

% \centering
% \begin{tikzpicture}[
%   node/.style={
%     rectangle, rounded corners=3pt, minimum width=3cm, minimum height=6mm,
%     draw=black, thick, fill=#1!20, align=center
%   },
%   labarrow/.style={midway, font=\scriptsize, fill=white, inner sep=1pt},
%   >=Latex
% ]

% % Nodo ra\'iz (siempre visible)
% \node[node=cyan] (n1) at (0,0) {Sistemas\\de transici\'on};

% \pause
% \node[node=green!60!black, right=2cm of n1] (n2) {Cadenas\\de Markov};

% % --- Paso 3: Nodo debajo de n2 ---
% \pause
% \node[node=yellow!60!black, below=1cm of n2] (n4) {Procesos de\\decisi\'on de Markov};
% \draw[->] (n2.south) -- node[labarrow, right] {agrego elecciones} (n4.north);
% \draw[->] (n1.south east) to node[labarrow] {agrego probabilidades} (n4.north west);

% % --- Paso 2: Nodo derecho ---
% \pause
% \only<4>{\node[node=orange!60!black, below=1cm of n1] (n3) {Juegos\\no-estoc\'asticos};}
% \draw[->] (n1) to node[labarrow, left] {agrego adversarialidad} (n3);

% % --- Paso 4: Nodo final centrado ---
% \pause
% \node[node=purple!60!black] at (2.75,-5) (n5) {Juegos\\estoc\'asticos};
% \draw[->] (n3.south) to node[labarrow] {agrego probabilidad} (n5);
% \pause
% \draw[->] (n4.south) to node[labarrow] {agrego adversarialidad} (n5);

% \end{tikzpicture}
% \end{frame}

% \begin{frame}{Formalicemos un poco...}
%     mdp

%     como analizamos a este modelo en el contexto de de verificacion: a traves de ver probabilidades de comportamientos (presentar cadenas de markov ? ¿

%     definir aca las estrategias o en objetivos? como se introducen?
% \end{frame}

% \begin{frame}{Ejemplo de proceso de decisi\'on de Markov}

% \end{frame}

\begin{frame}{Juego estoc\'astico}
	\pause
	% Un juego estocastico es en principio otro modelo matematico que se usa para representar sistemas en el ambito de la verificacion formal, y es de ellos que nace luego la idea de juego estocastico politopico. Entonces bueno matematicamente un juego estocastico es una tupla formada por 4 cositas, 
	\begin{block}{Juego estoc\'astico}
		Un juego estoc\'astico es una tupla $\G = (\St, (\St_\cuad, \St_\diam), \Act, \theta)$ donde:
		\begin{itemize}
			\item $\St$ es un conjunto finito de estados con $\St_\cuad, \St_\diam \subseteq \St$ siendo una partici\'on de \'el,
			\item $\Act$ es un conjuno finito de acciones, y
			\item $\theta : \St \times \Act \times \St \rightarrow [0,1]$ es una funci\'on de transici\'on probabil\'istica tal que para cada $s \in \St$, $\theta(s, a, \cdot) \in \dist(\St)$ o  $\theta(s, a, \St) = 0$.
		\end{itemize}
	\end{block}
\end{frame}

\begin{frame}{Juego estoc\'astico}

	primero estados que pertenecen a distintos jugadores -> que el estado
	pertenezca a un jugador significa que ese jugador elige la accion desde ahi
	luego transiciones luego acciones luego probabilidades

	elegir una accion determina una distribucion de probabilidad sobre en qu\'e
	estado estar\'a el juego a continuaci\'on
\end{frame}

\begin{frame}{Ejemplo de juego estoc\'astico}
	Tenemos dos robotcitos, generar imagenes de estos

	Decir c\'omo vamos a representar su estado

\end{frame}

\begin{frame}{Ejemplo de juego estoc\'astico}
	\begin{align*}
		&    \St_\cuad = \cuad \\
		&    \St_\diam = \diam \\
		&    \Act = \{\uparrow, \rightarrow, \downarrow, \leftarrow\} \\
	\end{align*}

\end{frame}

\begin{frame}{Juego estoc\'astico polit\'opico}
	\begin{block}{Juego estoc\'astico polit\'opico II}
		Un juego estoc\'astico polit\'opico es una extensi\'on de los juegos estoc\'asticos. \pause Difiere de estos porque los jugadores pueden elegir de entre un conjunto \textbf{infinito de acciones}. \pause Este conjunto es caracterizado a trav\'es de \textbf{desigualdades lineales} cuyo espacio de soluciones forma un \textbf{politopo}.
	\end{block}
\end{frame}

\begin{frame}{Politopos}
	\begin{block}{Politopo convexo}
		Un politopo convexo en $\mathbb{R}^n$ es un conjunto acotado $K = \{ x
			\in \mathbb{R}^n | Ax \leq b\}$ con $A \in \mathbb{R}^{m\times n}$ y $ b \in
			\mathbb{R}^m$, para alg\'un $m \in \mathbb{N}$.
	\end{block}
	\pause

	\begin{itemize}
		\item Las funciones en $\mathbb{R}^S$ se pueden pensar como vectores y podemos hablar
		      de politopos en $\mathbb{R}^S$. \pause
		\item El conjunto de las distribuciones de probabilidad en $S$ forma un politopo
		      convexo. \pause
		\item  Definimos como \( \DPoly(S) \) al conjunto de politopos convexos cuyos
		      elementos son funciones de probabilidad sobre $S$.
	\end{itemize}

	% Trabajando con PSGs lo que hacemos es pensar en politopos cuyos elementos (vectores x en las desigualdades Ax <= b) son funciones de probabilidad sobre $S$. Pensamos en restricciones que cumplen las funciones de probabilidad que pueden ser elegidas por los jugadores.
\end{frame}

\begin{frame}{Juego estoc\'astico polit\'opico}
	\begin{block}{Juego estoc\'astico polit\'opico}
		Un juego estoc\'astico polit\'opico (PSG, por sus siglas en ingl\'es) es una estructura \( \K = (\St, (\St_\square, \St_\diamondsuit), \Theta) \) tal que \( \St \) es un conjunto finito de estados particionado en $\St_\cuad$ y $\St_\diam$, (es decir, \( \St = \St_\square \biguplus \St_\diamondsuit \)) y \( \Theta : \St \to \mathcal{P}_f(\textsf{DPoly}(\St)) \).
	\end{block}
\end{frame}

\begin{frame}{Interpretaci\'on de un juego estoc\'astico polit\'opico}
	\begin{block}{Interpretaci\'on de un PSG}
		La interpretaci\'on del juego estoc\'astico polit\'opico \( \K \) se define por el juego estoc\'astico \( \Gk = (\St, (\St_\square, \St_\diamondsuit), \Act, \theta) \), donde \( \Act = \bigcup_{s \in \St} \Theta(s) \times \textsf{Dist}(\St) \) y
		\[
			\theta(s, (K, \mu), s') =
			\begin{cases}
				\mu(s') & \text{si } K \in \Theta(s) \text{ y } \mu \in K \\
				0 & \text{en otro caso}.
			\end{cases}
		\]
	\end{block}
\end{frame}

\begin{frame}{Ejemplo de juego estoc\'astico polit\'opico}

\end{frame}

\begin{frame}{Formalizando la noci\'on de elecci\'on}
	% Casi que entendimos la mitad del titulo, antes de pasar a la otra mitad quiero que retomemos un poquito la primera definicion que vimos de juego estocastico politopico. Dijimos que es un modelo \textbf{matematico} que nos servira para representar un sistema. Dijimos a su vez que en los sistemas que representamos tenemos dos jugadores o agentes que toman decisiones, pero 

	\begin{center}
		¿C\'omo definimos matem\'aticamente la noci\'on de elecci\'on de los jugadores?
		\pause

		A trav\'es de estrategias
	\end{center}
	\pause

	% Decir del truco. Yo cuando juego al truco si tengo un 4 y un 5 en la mano la mayoria de las veces no te voy a cantar truco, pero capaz dependiendo un poco de c\'omo viene la mano a veces si te puedo mentir y cantar truco. Te diria que si vengo perdiendo hay un 50 porciento de chances de que te cante o no y si vengo ganando nunca te voy a cantar con un 4 y un 5. Esa idea de formalizacion de lo que yo hago jugando al truco es un poco la misma idea de estrategia.

	% Si me escucharon habr\'an visto que mencione esta idea de dependiendo como viene la mano es que decido que hacer. asi que antes de definir formalmente lo que es una estrategia en realidad vamos a tener que formalizar esta idea de como viene la mano. 

	\begin{block}{Camino en un PSG}
		Un camino en un PSG es una secuencia infinita que alterna entre estados y acciones. Formalmente, un camino es un
		$$\omega = (s_0, (K_0, \mu_0), s_1, (K_1, \mu_1), \dots),$$
		tal que $s_i \in \St$, $(K_i, \mu_i) \in \Act(s_i)$ y $\mu_i(s_{i+1}) > 0$ para todo $i \geq 0$.

		\vspace{0.1cm}
		Denotaremos con $\Pathsfin^i$ al conjunto de los prefijos de camino
		finitos que terminan en un estado $s \in \St_i$.
	\end{block}

	% definiremos como comportamiento en un juego o camino a una secuencia alternante entre estados y acciones. llamaremos paths al conjunto de todos los posibles caminos en un juego y paths fin a todos los fragmentos finitos de caminos posibles en un juego.
\end{frame}

\begin{frame}{Estrategias}

	\only<1,2>{
		\begin{block}{Estrategia en un PSG}
			Una estrategia $\pi_i$ para el jugador $i$ ($i \in \{\cuad, \diam\}$) en un PSG ser\'a una funci\'on $\pi_{{i}} : \Pathsfin^{i} \rightarrow {\PMeas}({\Act})$ que asigna una {medida de probabilidad} sobre las {acciones} a cada ${\omega} {s}$ tal que $\pi_i(\omega s)(\Act(s)) = 1$.
		\end{block}}

	\only<3>{
		\begin{block}{Estrategia en un PSG}
			Una estrategia $\pi_i$ para el jugador $i$ ($i \in \{\cuad, \diam\}$) en un PSG ser\'a una funci\'on $\pi_{{i}} : \Pathsfin^{i} \rightarrow \textcolor{orange}{\PMeas}({\Act})$ que asigna una \textcolor{orange}{medida de probabilidad} sobre las {acciones} a cada ${\omega}{s}$ tal que $\pi_i(\omega s)(\Act(s)) = 1$.
		\end{block}}

	\only<4>{
		\begin{block}{Estrategia en un PSG}
			Una estrategia $\pi_i$ para el jugador $i$ ($i \in \{\cuad, \diam\}$) en un PSG ser\'a una funci\'on $\pi_{\textcolor{blue}{i}} : \Pathsfin^{i} \rightarrow \textcolor{orange}{\PMeas}({\Act})$ que asigna una \textcolor{orange}{medida de probabilidad} sobre las {acciones} a cada ${\omega} {s}$ tal que $\pi_i(\omega s)(\Act(s)) = 1$.
		\end{block}}

	\only<5>{
		\begin{block}{Estrategia en un PSG}
			Una estrategia $\pi_i$ para el jugador $i$ ($i \in \{\cuad, \diam\}$) en un PSG ser\'a una funci\'on $\pi_{\textcolor{blue}{i}} : \Pathsfin^{i} \rightarrow \textcolor{orange}{\PMeas}(\textcolor{magenta}{\Act})$ que asigna una \textcolor{orange}{medida de probabilidad} sobre las \textcolor{magenta}{acciones} a cada ${\omega}{s}$ tal que $\pi_i(\omega s)(\Act(s)) = 1$.
		\end{block}}

	\only<6>{
		\begin{block}{Estrategia en un PSG}
			Una estrategia $\pi_i$ para el jugador $i$ ($i \in \{\cuad, \diam\}$) en un PSG ser\'a una funci\'on $\pi_{\textcolor{blue}{i}} : \Pathsfin^{i} \rightarrow \textcolor{orange}{\PMeas}(\textcolor{magenta}{\Act})$ que asigna una \textcolor{orange}{medida de probabilidad} sobre las \textcolor{magenta}{acciones} a cada ${\omega} \textcolor{brown}{s}$ tal que $\pi_i(\omega s)(\Act(s)) = 1$.
		\end{block}}

	\only<7,8>{
		\begin{block}{Estrategia en un PSG}
			Una estrategia $\pi_i$ para el jugador $i$ ($i \in \{\cuad, \diam\}$) en un PSG ser\'a una funci\'on $\pi_{\textcolor{blue}{i}} : \Pathsfin^{i} \rightarrow \textcolor{orange}{\PMeas}(\textcolor{magenta}{\Act})$ que asigna una \textcolor{orange}{medida de probabilidad} sobre las \textcolor{magenta}{acciones} a cada $\textcolor{teal}{\omega} \textcolor{brown}{s}$ tal que $\pi_i(\omega s)(\Act(s)) = 1$.
		\end{block}}

	\pause
	La estrategia nos permite determinar \pause \textcolor{orange}{con qu\'e
		probabilidad} \pause \textcolor{blue}{un jugador} \pause elige
	\textcolor{magenta}{determinada acci\'on} \pause desde
	\textcolor{brown}{determinado estado} \pause (con alguna \textcolor{teal}{determinada
		historia}).

	\pause
	Luego, con la acci\'on se determina cu\'al ser\'a el pr\'oximo estado del
	juego.
\end{frame}

\begin{frame}{Algo m\'as sobre estrategias}
	Existen distintas maneras de clasificar estrategias:

	\begin{itemize}
		\item seg\'un si son deterministas o randomizadas.
		\item seg\'un si son sin memoria, con memoria finita o generales.
	\end{itemize}

\end{frame}

\subsection{¿Qu\'e son los objetivos de Rabin?}

\begin{frame}{Objetivos en un juego}
	Podemos definir a los objetivos en un juego como:
	\begin{itemize}
		\item Maneras de codificar la especificaci\'on deseada para nuestro sistema/modelo
		\item Conjuntos de comportamientos/caminos (secuencias de estados)
	\end{itemize}
	Ejemplos de objetivos podr\'ian ser:
	\begin{itemize}
		\item El robot X llega a la esquina inferior izquierda. \pause
		\item El robot X siempre pasa por la esquina superior izquierda. \pause
		\item El robot X solo pasa por la parte superior del tablero hasta llegar a la
		      esquina inferior derecha.
	\end{itemize}
\end{frame}

\begin{frame}{Objetivos de alcanzabilidad y seguridad}

\end{frame}

\begin{frame}{Objetivos de B\"uchi y co-B\"uchi}

\end{frame}

\begin{frame}{Objetivos de Rabin y Streett}

\end{frame}

\begin{frame}{¿Por qu\'e nos interesan los objetivos de Rabin?}

\end{frame}

\section{Aportes de la tesina}

\subsection{¿Cu\'al fue la pregunta de investigaci\'on?}

\begin{frame}{¿Cu\'al fue la pregunta de investigaci\'on?}
	\begin{center}
		\textquestiondown Desde qu\'e estados gana $\cuad$ en un juego estoc\'astico polit\'opico con objetivo de Rabin?

		\vspace{0.7cm}
		\pause
		\textquestiondown C\'omo puedo calcularlos?

		\vspace{0.7cm}
		\pause
		\textquestiondown C\'omo puedo construir una estrategia ganadora?
	\end{center}
\end{frame}

\subsection{¿C\'omo hicimos para responderla?}

\begin{frame}{La idea general}
	\textbf{Hicimos una reducci\'on a otra clase de juego}
	\pause
	\begin{enumerate}
		\item Presentamos la \textit{desrandomizaci\'on} de un juego estoc\'atico
		      polit\'opico. \pause
		\item Probamos que los conjuntos ganadores para $\cuad$ en ambos juegos son
		      equivalentes. \pause
	\end{enumerate}
	Con esto, adaptamos un \textbf{algoritmo} de ejecuci\'on simb\'olica para el c\'alculo de este conjunto, analizamos su \textbf{complejidad} y damos un m\'etodo para la obtenci\'on de \textbf{estratgias ganadoras}.
\end{frame}

\begin{frame}{Desrandomizaci\'on de un PSG}

	\begin{block}{Juego de dos jugadores con adversario justo}
		Sea $G = (V, (V_\cuad, V_\diam), E)$ un juego de grafo de dos jugadores y
		$E^{l} \subseteq (V_\diam \times V) \cap E$ un conjunto dado de aristas que
		deben ser tomadas infinitamente a menudo si el estado del que parten es
		visitado infinitamente a menudo.

		Denotamos con $G^{l} = \langle G, E^{l} \rangle$ a un juego de adversario
		justo.
	\end{block}

	% Lo importante para acordarse de esta definici\'on es que este es un juego
	% totalmente determinista es decir no hay probabilidades asociadas a ninguna de
	% las transiciones. Lo que si nos aseguramos de que en ciertos momentos al menos
	% el jugador diamante se comporta de manera justa. Eso es todo.

	% La idea ahora va a ser ver como se construye la desrandomizaci\'on a partir de
	% un juego estocastico politopico
\end{frame}

\begin{frame}{Desrandomizaci\'on de un PSG}
	Dada la interpretaci\'on de un PSG $\Gk = (\St, (\St_\cuad, \St_\diam), \Act,
		\theta)$, y denotando con $\Vs$ al conjunto $\{\supp(\alpha) | \alpha \in
		\Act(s)\}$ para cada $s \in \St$, se define la \textit{desrandomizaci\'on} de
	$\Gk$, $\DGk := (( V, V_\cuad, V_\diam, E ), E^l )$, como el juego de grafo de
	2 jugadores justo donde
	\begin{alignat*}{2}
		&\tilde V &&= \{v_{V} \mid \exists s \in \St \text{ tal que } V \in \Vs\}, \\
		&V &&= \St \cup \tilde V, \\
		&V_\cuad &&= \St_\cuad, \\
		&V_\diam &&= \St_\diam \cup \tilde V, \\
		&E^l &&= \{(v_{V}, s') \mid s' \in V\}, \text{   y } \\
		&E &&= \{(s, v_{V}) \mid V \in \Vs\} \cup E^l.\\
	\end{alignat*}

	Explicar la condici\'on de equidad en base a los nuevos vertices
\end{frame}

\begin{frame}{Ejemplo de desrandomizaci\'on de un PSG}
	repetir figura del psg para refrescar memoria
\end{frame}

\begin{frame}{Ejemplo de desrandomizaci\'on de un PSG}
	\begin{itemize}
		\item $V_{\substack{c(0,0) \\ \ (1,0)}} = \left\{\left\{\substack{d(0,1) \\ \ (1,0)}, \substack{d(1,1) \\ \ (1,0)}, \substack{d(1,0) \\ \ (1,0)}\right\}\right\}$
		\item $V_{\substack{d(0,1) \\ \ (1,0)}} = \left\{\left\{\substack{c(0,1) \\ \ (0,0)}\right\}, \left\{\substack{c(0,1) \\ \ (1,1)}\right\}\right\}$
		\item $V_{\substack{d(1,1) \\ \ (1,0)}} = \left\{\left\{\substack{c(1,1) \\ \ (0,0)}\right\}, \left\{\substack{c(1,1) \\ \ (1,1)}\right\}\right\}$
		\item $V_{\substack{d(1,0) \\ \ (1,0)}} = \left\{\left\{\substack{c(1,0) \\ \ (0,0)}\right\}, \left\{\substack{c(1,0) \\ \ (1,1)}\right\}\right\}$
	\end{itemize}
\end{frame}

\begin{frame}{Ejemplo de desrandomizaci\'on de un PSG}
	\begin{tikzpicture}[scale=0.55, ->]
		\tikzset{transform shape,              % important: scale affects node shapes & text
			x=1.4cm, y=1cm,            % compress horizontal / vertical distances
			state/.style={draw, circle, minimum size=1.2cm, font=\normalsize},
			square state/.style={draw, rectangle, minimum size=1.4cm, font=\normalsize, align=center},
			diamond state/.style={draw, diamond, aspect=1.5, minimum size=1.5cm, font=\normalsize, align=center},
			every node/.style={inner sep=4pt}
		}

		% Nodo ra\'iz a la izquierda
		\node[square state] (left) at (0, 0) {$(0,0)$\\$(1,0)$};

		% Nodo Vleft
		\node[diamond state] (Vleft) at (3,0) {$V^1_{\substack{c(0,0) \\ \ (1,0)}}$};

		% Nivel intermedio (rombos) centrados horizontalmente, espaciados verticalmente
		\node[diamond state] (d1) at (6, 4) {$(0,1)$\\$(1,0)$};
		\node[diamond state] (d2) at (6, 0) {$(1,1)$\\$(1,0)$};
		\node[diamond state] (d3) at (6, -4) {$(1,0)$\\$(1,0)$};

		% Nodos Vd
		\node[diamond state] (Vs1) at (9, 5) {$V^1_{\substack{d(0,1) \\ \ (1,0)}}$};
		\node[diamond state] (Vs2) at (9, 3) {$V^2_{\substack{d(0,1) \\ \ (1,0)}}$};

		\node[diamond state] (Vs3) at (9, 1.0) {$V^1_{\substack{d(1,1) \\ \ (1,0)}}$};
		\node[diamond state] (Vs4) at (9, -1.0) {$V^2_{\substack{d(1,1) \\ \ (1,0)}}$};

		\node[diamond state] (Vs5) at (9, -3) {$V^1_{\substack{d(1,0) \\ \ (1,0)}}$};
		\node[diamond state] (Vs6) at (9, -5) {$V^2_{\substack{d(1,0) \\ \ (1,0)}}$};

		% Nodos hojas (cuadrados), suficientemente separados
		\node[square state] (s1) at (12, 5) {$(0,1)$\\$(0,0)$};
		\node[square state] (s2) at (12, 3) {$(0,1)$\\$(1,1)$};

		\node[square state] (s3) at (12, 1.0) {$(1,1)$\\$(0,0)$};
		\node[square state] (s4) at (12, -1.0) {$(1,1)$\\$(1,1)$};

		\node[square state] (s5) at (12, -3) {$(1,0)$\\$(0,0)$};
		\node[square state] (s6) at (12, -5) {$(1,0)$\\$(1,1)$};

		% Arista de left a Vleft
		\draw[-> ] (left) to (Vleft);

		% Aristas desde left a cada rombo
		\draw[-> ] (Vleft) to (d1);
		\draw[-> ]  (Vleft) to (d2);
		\draw[-> ] (Vleft) to (d3);

		% Aristas desde cada rombo a dos Vs
		\draw[-> ] (d1) to (Vs1);
		\draw[-> ] (d1) to (Vs2);

		\draw[-> ] (d2) to (Vs3);
		\draw[-> ] (d2) to (Vs4);

		\draw[-> ] (d3) to (Vs5);
		\draw[-> ] (d3) to (Vs6);

		% Aristas desde cada Vs a cuads
		\draw[-> ] (Vs1) to (s1);
		\draw[-> ] (Vs2) to (s2);

		\draw[-> ] (Vs3) to (s3);
		\draw[-> ] (Vs4) to (s4);

		\draw[-> ] (Vs5) to (s5);
		\draw[-> ] (Vs6) to (s6);

		% Puntos suspensivos para indicar que contin\'ua
		\node at (13,5) {$\cdots$};
		\node at (13,3) {$\cdots$};

		\node at (13,1) {$\cdots$};
		\node at (13,-1) {$\cdots$};

		\node at (13, -3) {$\cdots$};
		\node at (13, -5) {$\cdots$};

	\end{tikzpicture}
\end{frame}

\begin{frame}{La prueba de igualdad de conjuntos ganadores}

	\begin{block}{Teorema principal}
		Sean:
		\begin{itemize}
			\item $\W$: conjunto de estados desde los que gana $\cuad$ en $\DGk$, y
			\item $\W^{as}$: conjunto de estados desde los cuales $\cuad$ gana con probabilidad 1 en $\Gk$ \tiny{con alguna estrategia de memoria finita}. \normalsize
		\end{itemize}

		Entonces: $$\W =\W^{as}$$ \vspace{0.1cm}
	\end{block}
	%     \begin{block}{Teorema principal}
	% 	\label{teocuali}
	% 	Sea $\Gk = (\St, (\St_\cuad, \St_\diam), \Act, \theta)$ la interpretaci\'on de un juego estoc\'astico polit\'opico, $R = \{(E_1, F_1), \dots, (E_d, F_d)\}$ una condici\'on de Rabin sobre $\St$, con su especificaci\'on LTL $\varphi$ y sea $\DGk$ su desrandomizaci\'on.

	% 	Sea $\W \subseteq \St$ el conjunto de todos los estados desde los cuales el
	% 	jugador $\cuad$ gana en $\DGk$ y sea $\W^{as}$ el conjunto de v\'ertices desde
	% 	los cuales el jugador $\cuad$ gana con probabilidad 1 en $\Gk$ con alguna
	% 	estrategia de memoria finita. Entonces, $\W = \W^{as}$.

	% 	Es m\'as, a partir de una estrategia ganadora en $\DGk$ se puede construir una
	% 	estrategia ganadora en $\Gk$, y viceversa.
	% \end{block}
\end{frame}

\begin{frame}{La prueba de igualdad de conjuntos}
	Probamos la doble contenci\'on:
	\begin{enumerate}
		\item $\W \subseteq \W^{as}$:

		      \begin{itemize}
			      \item Sea $s \in \W$. Entonces, existe una estrategia ganadora $\sigma_\cuad$ en
			            $\DGk$ desde $s$.
			      \item Construimos una estrategia $\pi_\cuad$ en $\Gk$ a partir de $\sigma_\cuad$.
			      \item Como $\pi_\cuad$ está definida a partir de $\sigma_\cuad$,
			            $\Pr^{\pi_\cuad}_s(\varphi^l \rightarrow \varphi) = 1$, donde $\varphi^l$ es la
			            suposici\'on de justicia para el juego desrandomizado y $\varphi$ es la
			            condici\'on de victoria de Rabin.
			      \item Luego, probamos que $\ProbG(\neg \varphi^l) = 0$.
			      \item Por lo tanto, $\Pr^{\pi_\cuad}_s(\varphi) = 1$ y $s \in \W^{as}$.
		      \end{itemize}
		      \seti
	\end{enumerate}
\end{frame}

\begin{frame}{La prueba de igualdad de conjuntos}
	Probamos la doble contenci\'on:
	\begin{enumerate}
		\conti
		\item $\W^{as} \subseteq \W$:

		      \begin{itemize}
			      \item Sea $s \in \W^{as}$. Entonces, existe una estrategia ganadora $\pi_\cuad$ en
			            $\Gk$ desde $s$.
			      \item Construimos una estrategia $\sigma_\cuad$ en $\DGk$ a partir de $\pi_\cuad$.
			      \item Probamos que para todo camino $\rho$ desde $s$ generado por $\sigma_\cuad$ y
			            alguna estrategia $\sigma_\diam$ en $\DGk$, se cumple que $\Inf_C(\rho) =
				            \Inf(\rho) \cap \St$ cumple con la especificaci\'on de Rabin.
			      \item Por lo tanto, $s \in \W$.
		      \end{itemize}
	\end{enumerate}
\end{frame}

\begin{frame}{Procesos de Decisi\'on de Markov Polit\'opicos}

\end{frame}

\subsection{¿Qu\'e resultados pr\'acticos obtuvimos a partir de ella?}

\begin{frame}{Algoritmo para el c\'alculo de regiones ganadoras}
	Adaptamos un algoritmo planteado por Banerjee et al. que:
	\begin{itemize}
		\item calcula el conjunto de estados ganadores en un juego de grafo de dos jugadores
		      con adversario justo y objetivo de Rabin, \pause
		\item está planteado en $\mu$-c\'alculo y caracterizado mediante una expresión de
		      punto fijo, y \pause
		\item es de ejecuci\'on simb\'olica. Es decir, trabaja usando representaciones
		      compactas de conjuntos de estados y aplica transformadores sobre ellos.
	\end{itemize}
	% Agregar slide mostrando el algoritmo enunciandolo bien
\end{frame}

\begin{frame}{Sint\'esis de estrategia ganadora en un PSG}

\end{frame}

\section{Conclusi\'on}

\subsection{Resumen de los aportes}

\begin{frame}{Resumen de los aportes}
	\pause
	\begin{enumerate}
		\item Presentamos la \textit{desrandomizaci\'on} de un PSG. \pause
		\item Establecimos una correspondencia entre caminos y estrategias en un PSG y en su
		      desrandomizaci\'on, con lo que pudimos demostrar la igualdad entre los
		      conjuntos ganadores en ambos. \pause
		\item Y con eso obtuvimos:
		      \begin{itemize}
			      \pause
			      \item un proceso para sintetizar estrategias ganadoras en un PSG con objetivo de
			            Rabin, y \pause
			      \item un algoritmo para el c\'alculo de conjunto de estados ganadores en un PSG con
			            objetivo de Rabin.
		      \end{itemize}
		\item Además, introdujimos el concepto de procesos de decisión de Markov politópicos
		      y probamos unos teoremas básicos sobre ellos.
	\end{enumerate}
\end{frame}

\subsection{Trabajo futuro}

\begin{frame}{Trabajo futuro}
	%podria decir que posible trabajo futuro 
	La idea principal ser\'ia (fue) intentar responder: ¿cu\'al es la probabilidad
	m\'axima que tiene un jugador de ganar desde un estado?

	Ideas para responder esa pregunta:
	\begin{itemize}
		\item mostrar que si una familia de estrategias es suficiente para el c\'alculo de
		      regiones ganadoras entonces tambi\'en es suficiente para el c\'alculo de
		      valores de los estados.
		\item ver que el encontrar el valor en un estado de un PSG para un objetivo de Rabin
		      puede ser equivalente a buscar el valor de un estado en la interpretaci\'on
		      extrema del juego.
	\end{itemize}
\end{frame}

\begin{noheadline}
	\begin{frame}{}
		\begin{center}
			¡Muchas gracias!
		\end{center}
	\end{frame}
\end{noheadline}

\begin{noheadline}
	\begin{frame}<beamer>[noframenumbering]
		Idea escrita

		Pasaje a imperativo
	\end{frame}
\end{noheadline}

\end{document}
